% !TEX root = ../MasterThesis.tex
\normallinespacing

\chapter{Introduction}

A stochastic system that must not fail cannot be trusted on the strength of a few successful
simulations. Simulation samples behaviour; it never rules out the rare
trajectory that ends in disaster. Examples of risky systems include a controller keeping a drone aloft, a policy managing financial risk, or a robot navigating on noisy sensors. What such systems demand instead is a
\emph{formal guarantee}: a proof that the system meets its required specification with
quantified confidence, holding over every admissible initial condition and every
realisation of the noise. This thesis is about producing such guarantees
automatically, and most importantly about producing them
\emph{quickly}.

\section{Motivation}

The dominant route to a formal guarantee is a \emph{certificate}: a single
function whose existence witnesses a proof that a system satisfies a property,
reducing an intractable question about infinitely many trajectories to a local,
checkable inequality on one function~\cite{lyapunov1992general}. Verifying a candidate then
amounts to establishing one inequality everywhere on the domain.

Classical synthesis fixes the certificate's functional form in advance (for example, a
polynomial for sum-of-squares programming~\cite{parrilo2000}) buying soundness at the price of expressiveness, and scaling
poorly beyond low dimensions and simple dynamics. The modern alternative is to
\emph{learn} the certificate as a neural network, expressive enough for
high-dimensional, non-polynomial, and even learned dynamics, but forfeiting any
closed-form guarantee: fitting a network on finitely many sampled states says
nothing about the uncountably many states between them.
\emph{Counterexample-guided inductive synthesis} (CEGIS) recovers soundness by
interleaving a \emph{learner}, which trains the network on a dataset of states,
with a \emph{verifier}, which checks the drift condition over the whole domain
and, on failure, returns a counterexample for the learner to
repair~\cite{chatterjee2023learnerverifier, abate2021fossil}. The loop repeats
until the verifier certifies the candidate.

A verifier inside CEGIS is doing two different jobs: while the candidate is invalid, the loop needs only a
\emph{counterexample} to train on; only once the candidate is finally valid
does it need a \emph{proof} that no counterexample exists. Despite this, the prevailing practice
discharges both jobs with a \emph{single} method running
every round, doubling as the counterexample generator on the rounds it
rejects. Sound verification is expensive because it must reason over
the entire domain. Paying that cost every round, when most rounds need
only a counterexample, is the bottleneck this thesis sets out to remove.

\section{Objectives}

The guiding hypothesis, stated formally in
Chapter~\ref{sec:preliminaries}, is that giving each of these tasks its own
dedicated algorithm, a cheap refuter to hunt counterexamples and a sound
verifier reserved for the final proof, yields a substantial speed-up.
We provide multiple contributions to establish and support this
claim.\footnote{An open-source implementation of every verifier,
refuter, experiment driver, and figure in this thesis is available at
\url{https://github.com/will-bailkoski/streamlining-neural-certificates}.}

\begin{itemize}
    \item \textbf{A unified survey and implementation of certificate
      verification for stochastic systems.} We bring three verifier families
      --- symbolic solvers (SMT, MILP), abstract interpretation (the LiRPA
      family), and sound sampling methods --- under a single interface, and map
      each onto a benchmark suite. The suite spans dynamics classes, noise models, and dimensions,
      highlighting where each method succeeds and fails.

    \item \textbf{A dedicated treatment of refutation as black-box
      optimisation.} We cast counterexample search as global optimisation and
      benchmark six optimisers. Using uninformed,
      Lipschitz-aware, and population-based search algorithms, we test on a standard benchmark, then on a library of real \emph{faux
      certificates} harvested from genuine CEGIS runs, finding that a well-tuned
      refuter matches the verifier's counterexample quality at one to three
      orders of magnitude lower cost.

    \item \textbf{Empirical evidence that separating refutation from verification
      streamlines synthesis.} Placing a fast refuter pre-screen in front of a
      sound verifier reduces CEGIS to a sequence of cheap search-and-retrain
      rounds followed by a single certification call. The speed-up scales with
      the cost of that call, improving the hardest benchmark from a multi-hour run to minutes.

    \item \textbf{Tighter Lipschitz constants for sample-based verification.}
      Because verification cost grows as a power of the drift-Lipschitz constant,
      we tighten estimates on the object sample-based verifiers consume: an exact drift-Lipschitz
      constant for linear systems, computed directly rather than by loose
      composition, and sound per-axis (anisotropic) constants. The two savings
      multiply, reducing the number of samples needed to cover a continuous
      domain by more than an order of magnitude in four dimensions.
\end{itemize}

\section{Structure of the Report}

Chapter~\ref{sec:preliminaries} defines the problem: the class of stochastic
systems, the ranking supermartingale certificate and its drift condition, the
CEGIS loop that synthesises it, and the central hypothesis the rest of the thesis
develops. For \textbf{Advanced Topics in AI}, Chapter~\ref{sec:lipschitz} is a focused study of the Lipschitz
constant and contributes two methods for tightening the constant for better verification performance. Chapter~\ref{chap:benchmarks} presents the benchmark suite, selecting
systems from the control and formal-methods literature so that difficulty rises
along interpretable axes. The two halves of the loop are then treated in
isolation: Chapter~\ref{chap:verification} surveys and evaluates the sound
verifiers, establishing the capability boundary and the cost of a sound call,
while Chapter~\ref{chap:refutation} develops refutation as cheap counterexample
search and benchmarks the refuters against it. Finally,
Chapter~\ref{chap:streamlining} assembles the two into the full pipeline and
tests the hypothesis directly, measuring the end-to-end speed-up of a refuter
pre-screen over a verifier-only baseline.


\newpage
